\section{Three positive residuals cannot include two units}
\label{sec:three-positive-two-units}

All far sets, protection relations, convenience and failure sets are
defined in the original graph.
For a far target $y$ of $x$, protection means
$y\to x$ and $N^-(y)\subseteq N^-(x)\setminus\{y\}$.
Write $\mathcal F(x)$ for the unprotected far targets and
$t(x)=|\mathcal F(x)|$.
As before, good missing pairs are oriented canonically: use the only
convenient direction, or use a fixed linear extension of the protection
partial order when both directions are convenient.
Protected-target compatibility ensures that an added convenient
$z\to a$ cannot be followed by an added arc from $a$ to a protected
far target of $z$.

\subsection{A simultaneous single-blame completion}

\begin{lemma}\label{lem:single-blame-unprotected-transfer}
Suppose every bad missing pair of an oriented graph $D$ can be oriented
so that all failure sources of its selected direction lie in $\{u\}$.
Orient all bad pairs this way and all good pairs canonically, obtaining
a tournament $T$.
For every $z\ne u$, put
$K_z=N_T^+(z)\setminus N_D^+(z)$.
Then
\[
 {
 N_T^{++}(z)\subseteq
 (N_D^{++}(z)\setminus K_z)
 \cup(\mathcal F(z)\setminus K_z),}
 \tag{TU1}\label{eq:single-blame-envelope}
\]
and
\[
 g_T(z)\ge g_D(z)+2|K_z|-t(z).
\]
If $K_z=\varnothing$, the second neighborhood is exactly unchanged.
\end{lemma}
\begin{proof}
All added first neighbors are missing partners, and therefore lie in
the disjoint union $N_D^{++}(z)\dot\cup\mathcal F(z)$.
We check that no protected far target $y$ of $z$ becomes a second
neighbor in $T$.

If the first intermediate vertex is an original outneighbor, an
original or convenient second arc cannot introduce an original far
head.  A selected bad second arc fails only at $u$, not at $z$.
If the first arc $z\to a$ is added and good, protection excludes an
original $a\to y$, and canonical compatibility chooses $y\to a$
whenever $ay$ is missing.

Now let $z\to a$ be an added bad arc.
Its opposite direction has a failure source $r\ne u$:
$r\to a$ and $z\in U_D(r)$.
The source cannot be $u$, since it already fails the selected direction.
Protection makes $y\in U_D(r)$.
An original $a\to y$ would force $a\to z$, contrary to missing $za$.
If $ay$ is missing, the direction $a\to y$ fails at $r$.
It cannot be a selected convenient direction when $ay$ is good,
or a selected bad direction when $ay$ is bad, because $r\ne u$.
Thus the completion points from $y$ to $a$.
This explicitly rules out paths using two added bad arcs.

Every genuinely new far second head is therefore unprotected.
Exclude all of $K_z$, now first neighbors, to obtain
\eqref{eq:single-blame-envelope}.
The degree inequality follows by removing $|K_z|$ from the combined
cardinality and adding $|K_z|$ to the first degree.
If $K_z$ is empty, every first intermediate is original and the first
case above prevents every new far head, whether protected or not.
Original second-neighborhood paths remain, so the second neighborhood
is unchanged.
\end{proof}

\begin{corollary}\label{cor:single-blame-unit-others}
Under the preceding orientation hypothesis, if $t(z)\le1$ for every
$z\ne u$, then $D$ has a Seymour vertex.
\end{corollary}
\begin{proof}
Suppose all vertices are deficient.
For $K_z=\varnothing$, the completed degree gap is $g_D(z)>0$.
For $|K_z|\ge1$, the bound gives
$g_T(z)\ge g_D(z)+2|K_z|-1>0$.
Thus every vertex except possibly $u$ is deficient in the same $T$.
All original outdegrees are positive; hence $T$ has no sink.
The published theorem giving at least two Seymour vertices in a
sink-free tournament \cite[Theorem~2]{HavetThomasse2000}
contradicts this simultaneous conclusion.
\end{proof}

\subsection{A unit--unit witness pair is itself bad}

The remainder assumes that $D$ is strongly connected, all-deficient
and minimal under deletion of arbitrary nonempty sets of arcs.
No restriction on the order or minimum outdegree is imposed.
The established residual and unit-target estimates are
\[
 \eta(x)=2\mu(x)+g(x)-1-|P(x)|\ge0,\qquad
 \eta(x)=1\Longrightarrow t(x)\le1.
\]

\begin{lemma}[The packet of an active unit source]
\label{lem:active-unit-full-boundary}
If $\eta(x)=1$ and $t(x)=1$, then $P(x)\ne\varnothing$ and
\[
 \partial P(x)=N_D^{--}(x).
\]
\end{lemma}
\begin{proof}
The sole unprotected far target prevents $x$ from being whole:
all far targets of a whole source are protected.
Hence $\mu(x)\ge1$, and $P(x)=\varnothing$ would give
$1=2\mu(x)+g(x)-1\ge2$.
So $P(x)$ is nonempty.
Set $C=\partial P(x)$, $Q=N_D^{--}(x)$,
$j=|Q\setminus C|$ and
$s=|N_D^+(x)\cap(Q\setminus C)|$.
The established target-bundle bounds give
\[
 s=0\Rightarrow\eta(x)\ge2j,\qquad
 s>0\Rightarrow\eta(x)\ge g(x)+2(j-s).
\]
If $s>0$, unit residual forces $g(x)=1$ and $j=s$.
The equal-defect unprotected-target bound would then give
$t(x)\le\eta(x)-g(x)=0$, a contradiction.
Thus $s=0$, and $1\ge2j$ forces $j=0$.
\end{proof}

\begin{lemma}\label{lem:unit-unit-witness-tail-bad}
Suppose opposite failure witnesses of a bad missing pair are vertices
$v,w$ with $\eta(v)=\eta(w)=1$.
Then $vw$ is itself a bad missing pair.
\end{lemma}
\begin{proof}
Label the original bad pair $ab$ so that
$v\to a$, $b\in U_D(v)$, $w\to b$ and $a\in U_D(w)$.
The witness pair $vw$ is missing.
Cross-witness protection shows
$\mathcal F(v)=\{b\}$ and $\mathcal F(w)=\{a\}$.
The heads $b,a$ are different from $w,v$, respectively, since
$w\to b$ and $v\to a$ are not loops.

A missing partner which is far is necessarily unprotected.
Thus the missing pair $vw$ is reachable in both directions in
exactly two steps: it uses neither source's sole unprotected head.
Lemma~\ref{lem:active-unit-full-boundary} gives
$w\in N_D^{--}(v)=\partial P(v)$ and
$v\in N_D^{--}(w)=\partial P(w)$.
Choose $p\in P(v)$ with $p\to w$, and $q\in P(w)$ with $q\to v$.
Then $v$ is far from $p$ and $w$ is far from $q$.
The direction $w\to v$ fails at $p$, and $v\to w$ fails at $q$.
Both directions fail, so $vw$ is bad.
\end{proof}

\begin{theorem}\label{thm:three-positive-two-units}
An arc-set-minimal strongly connected counterexample cannot have
positive residual support $\{u,v,w\}$ with
$\eta(v)=\eta(w)=1$, regardless of the positive value $\eta(u)$.
Equivalently, no allocation $(m,1,1)$ with $m\ge1$ is possible.
In particular the residual-five allocation $(3,1,1)$ is excluded.
\end{theorem}
\begin{proof}
Every directional failure source of a bad missing pair has positive
residual, by the zero-residual protection theorem.
Since $D$ is all-deficient, Ghazal's theorem
\cite[Theorem~2]{Ghazal2012} guarantees at least one bad pair.

The vertices $v,w$ cannot be opposite witnesses of any bad pair.
Otherwise Lemma~\ref{lem:unit-unit-witness-tail-bad} would make $vw$
bad as well.  Its two opposite failure witnesses would have to be
distinct positive-residual vertices outside $\{v,w\}$.
Only $u$ is available, a contradiction.

Consequently, for each bad pair, its two nonempty directional failure
sets cannot have a unit source on each side.
They are disjoint subsets of $\{u,v,w\}$.
One direction therefore has no source in $\{v,w\}$ and has precisely
$u$ as its failure source.
Choose this direction for every bad pair.
In cross-witness graph language, the graph has no edge $vw$ and
all its nonempty edges lie in the star centered at $u$.
This is an explicit choice of a single-blame orientation; no assertion
about arbitrary acyclic dependency digraphs is used.

The unit estimate gives $t(v),t(w)\le1$.
All other vertices except $u$ have residual zero, hence $t=0$.
Corollary~\ref{cor:single-blame-unit-others} now produces a Seymour
vertex of $D$, the required contradiction.
\end{proof}
